\section{How many rounds can there be?}
\label{sec:upper}

Here is the main theorem. Recall that rounds are counted with~\Hs1 as
the first one (Section~\ref{sec:counting}), and let
\[
   Q_n=\bigfloor{\frac{(n+1)^2}{4}}=\bigfloor{\frac{n+1}2}\bigceil{\frac{n+1}2},
\]
the quarter-squares $(Q_1,Q_2,Q_3,\ldots)=(1,2,4,6,9,12,16,20,25,30,36,42,\ldots)$.

\begin{lthm}{R}
Let $s\ge0$. Every ordered bipartite graph whose maximum matching has
size~$s$ makes Algorithm~H perform at most
\[
   R(s)=\bigfloor{\sqrt{4s+1}}-1
\]
rounds, and some ordered graph with $M=N=s$ attains this number.
Equivalently, $n$ rounds are possible if and only if $s\ge Q_n$.
\end{lthm}

The two formulations agree: $Q_n\le s$ means $(n+1)^2<4s+4$, that is
$(n+1)^2\le4s+3$; since squares are $0$ or $1$ modulo~$4$, this is the
same as $(n+1)^2\le4s+1$, or $n\le\floor{\sqrt{4s+1}}-1$.

This section proves the upper bound. Section~\ref{sec:staircase}
constructs graphs that attain it.

\subsection{A rectangle that must fit}

The argument is the one in the proof of Theorem~H, applied at the start
of a round instead of at the end.

\begin{lemma}
\label{lem:rectangle}
Suppose that before some round the current matching has $s'$ couples,
and that its shortest augmenting paths have level~$L$. Then
\[
   (s-s')\,L\ \le\ s .
\]
\end{lemma}

\begin{proof}
The union of the current matching with a maximum matching contains
$s-s'$ vertex-disjoint augmenting paths. Each has level at least $L$,
hence contains at least $L$ edges of the maximum matching. These edges
are distinct, and there are only $s$ of them.
\end{proof}

When $s=M=N$, the deficiency $s-s'$ is the number $f$ of free girls, and
the lemma says that the $f$ basic paths, drawn as rows of at least~$L$
cells each, contain an $f\times L$ rectangle inside the set of all $s$
cells (Fig.~\ref{fig:hyperbola}, left). Nothing requires the rows to have
equal numbers of cells; a basic path may have a much higher level than a
SAP.

\begin{figure}[ht]
\centering
% Left: basic paths as rows (s = 20); right: the (L,f) plane for s = 16.
\begin{tikzpicture}[x=.42cm,y=.42cm,baseline=(current bounding box.center)]
  % rows of 6,5,4,4 cells (= 19 cells) plus one cell on no basic path
  \foreach \len [count=\r] in {6,5,4,4}{
    \foreach \c in {1,...,\len}{
      \ifnum\c<5
        \draw[fill=sapcol!30,draw=black!35] (\c-1,-\r) rectangle (\c,-\r+1);
      \else
        \draw[fill=boyfill,draw=black!35] (\c-1,-\r) rectangle (\c,-\r+1);
      \fi}}
  \draw[fill=black!12,draw=black!35] (0,-5.15) rectangle (1,-4.15);
  \draw[decorate,decoration={brace,amplitude=4pt}] (-.25,-4)--(-.25,0)
    node[midway,left=4pt,font=\footnotesize] {$f$};
  \draw[decorate,decoration={brace,amplitude=4pt}] (0,.25)--(4,.25)
    node[midway,above=4pt,font=\footnotesize] {$L$};
\end{tikzpicture}\hspace{12mm}%
\begin{tikzpicture}[baseline=(current bounding box.center)]
\begin{axis}[
  width=8.4cm,height=5.8cm, xmin=1.5,xmax=7.6,ymin=0,ymax=9.6,
  xtick={2,...,7},ytick={1,...,9},
  xlabel={level $L$}, ylabel={free girls $f$},
  axis lines=left, axis line style={-}, tick label style={font=\footnotesize},
  label style={font=\small},
  legend style={draw=none,fill=none,at={(1.0,1.0)},anchor=north east,font=\footnotesize,
    cells={anchor=west}}]
  \addplot[forget plot,fill=boyfill,draw=none,domain=1.67:7.6,samples=80] {16/x} \closedcycle;
  \addplot[forget plot,domain=1.67:7.6,samples=80,black!55,thick] {16/x};
  \addplot[chcol,mark=*,mark size=1.8pt,thick,densely dashed]
     coordinates {(2,8) (3,5) (4,3) (5,2) (6,1)};
  \addlegendentry{Tchoukaillon schedule, $s=16$}
  \addplot[qcol,mark=square*,mark size=1.6pt,thick]
     coordinates {(2,6) (3,5) (4,4) (5,3) (6,2) (7,1)};
  \addlegendentry{staircase graph $\Gamma_7$, $s=16$}
  \node[font=\footnotesize,black!60,anchor=south west] at (axis cs:6.45,2.5) {$fL=16$};
\end{axis}
\end{tikzpicture}

\caption{Left: an example with $s=20$ and $f=4$. The basic paths of the
matching are drawn one per row, each with at least $L=4$ cells (one cell
per pair of the perfect matching); the grey cell lies on no basic path.
The shaded $f\times L$ rectangle shows why $fL\le s$. Right: every round
of every run is a point $(L,f)$ on or below the hyperbola $fL=s$. Both
trajectories shown take $s=16$; the staircase graph of
Section~\ref{sec:staircase} starts with fewer free girls but loses only
one per round, and so lasts longer.}
\label{fig:hyperbola}
\end{figure}

\subsection{Proof of the upper bound}

Suppose a run has $n\ge2$ rounds, so that there are $n-1$ successful
rounds after~\Hs1. Number these $i=1,2,\ldots,n-1$ and let round~$i$
have level~$L_i$ and start with deficiency $d_i=s-s'$.

\begin{itemize}
\item $L_i\ge i+1$. Indeed $L_1\ge2$ because the matching produced by~\Hs1 is
maximal, and the levels strictly increase from round to round.
\item $d_i\ge n-i$. Indeed each of the rounds~$i,i+1,\ldots,n-1$ adds at
least one couple.
\end{itemize}
By Lemma~\ref{lem:rectangle}, $s\ge(i+1)(n-i)$ for every $i$. The two
factors add up to $n+1$, so the product is largest when they are as
nearly equal as possible: taking $i+1=\ceil{(n+1)/2}$ gives
$s\ge\floor{(n+1)/2}\ceil{(n+1)/2}=Q_n$. A run with a single round needs
an edge, so $s\ge1=Q_1$. This proves the upper bound in Theorem~R.
\qed

\medskip
Hopcroft and Karp bounded the number of distinct path lengths, and
hence the number of phases, by $2\floor{\sqrt s}+2$ \cite[Theorem~3]{HK},
by essentially this argument applied at a single well-chosen phase.
Using it at every round gives the exact threshold~$Q_n$. What is new in
Theorem~R is that this threshold is attained.

\subsection{What a run with the most rounds must look like}
\label{sec:equality}

If $s=Q_n$ and a run has $n$ rounds, then the inequalities in the proof
are tight at a middle round~$i$, where $(i+1)(n-i)=Q_n$. That round has
level exactly $L=i+1$ and deficiency exactly $f=n-i$, and the $f$
disjoint augmenting paths supplied by a maximum matching are all SAPs,
which together use every one of its $s$ edges. So $\nu=f$ there, and in
the picture the $f\times L$ rectangle fills all $s$ cells, as in Knuth's
matrices. From that round on, every round removes exactly one SAP.

The rounds before the middle are where a worst case can differ from a
Tchoukaillon graph. There $f$ can exceed $L+1$, and then, if $\nu$ were
equal to $f$, Lemma~\ref{lem:15ab}(b) would force at least two SAPs in the
round. The
staircase graphs of Section~\ref{sec:staircase} avoid this by making some
basic paths longer than~$L$: those paths are counted in~$f$ but are not
SAPs, so $\nu\le L+1$ and a single SAP can meet all the shortest ones.
Fig.~\ref{fig:hyperbola} (right) shows the effect: the Tchoukaillon
schedule starts at the top of the hyperbola and must shed free girls
quickly, while the staircase graph starts lower and slides down along it,
one free girl per round.
